% -------------------------------------------------------------
%  RESUMOS
\setlength{\absparsep}{18pt} % ajusta o espaçamento dos parágrafos do resumo
% -------------------------------------------------------------
% resumo em PORTUGUÊS (OBRIGATÓRIO)
\begin{resumo}[Resumo]
 \begin{otherlanguage*}{brazilian}
    Este trabalho aborda o problema da estimação de efeitos causais médios em trajetórias temporais (FATE) no contexto de dados funcionais. Tradicionalmente, a literatura de inferência causal funcional recorre a métodos de suavização globais, como \emph{splines}, que impõem suposições de suavidade e regularidade frequentemente violadas na prática, além de dificultarem o processamento de dados em tempo real. Como alternativa, propomos o uso de transformadas de ondaletas implementadas via \emph{lifting scheme} como etapa de pré-processamento. Desenvolvemos o pacote \texttt{rLifting} em linguagem \texttt{R}, integrando rotinas em \texttt{C++} para garantir alta eficiência computacional e suporte a janelas móveis operacionalmente causais. Avaliamos o desempenho da proposta por meio de estudos de simulação de Monte Carlo, comparando estimadores \emph{offline} e \emph{online} com abordagens baseadas em \emph{splines}. Os resultados demonstram que o uso de ondaletas oferece maior precisão e robustez em cenários com descontinuidades e picos de frequência, mantendo coberturas estatísticas nominais e latência computacional reduzida. O trabalho conclui que a metodologia proposta é uma alternativa viável e superior para a inferência causal funcional em fluxos contínuos de dados.

    \textbf{Palavras-chave}: Inferência Causal. Dados Funcionais. Ondaletas. Lifting Scheme. Processamento em tempo real.
 \end{otherlanguage*}
\end{resumo}
% -------------------------------------------------------------
% resumo em INGLÊS (OBRIGATÓRIO)
\begin{resumo}[Abstract]
 \begin{otherlanguage*}{english}
    This work addresses the estimation of Functional Average Treatment Effects (FATE) within the context of functional data analysis. Conventional functional causal inference often relies on global smoothing methods, such as splines, which impose smoothness and regularity assumptions frequently violated in practice and hinder real-time data processing. As an alternative, we propose the use of wavelet transforms implemented via the lifting scheme as a pre-processing step. We developed the \texttt{rLifting} package in \texttt{R}, integrating \texttt{C++} routines to ensure high computational efficiency and support for operationally causal moving windows. We evaluated the proposal's performance through Monte Carlo simulation studies, comparing offline and online estimators with spline-based approaches. The results demonstrate that wavelets provide greater precision and robustness in scenarios with discontinuities and frequency spikes, maintaining nominal statistical coverage and low computational latency. The work concludes that the proposed methodology is a viable and superior alternative for functional causal inference in continuous data streams.
    
    \textbf{Keywords}: Causal Inference. Functional Data Analysis. Wavelets. Lifting Scheme. Real-time processing.
 \end{otherlanguage*}
\end{resumo}
% -------------------------------------------------------------
